\begin{abstract}
Adaptive retrieval requires a model to decide when its internal knowledge is
insufficient. Token entropy is a natural training-free signal, but it mixes
competition among plausible leading tokens with probability mass dispersed
across the rest of the vocabulary. We introduce \tailleak{}, a confidence score
that retains only the latter effect. At each generated position, it measures
the negative log probability mass retained by the model's top-\(k\) tokens and
sums this quantity over the answer. The score is available from one greedy
decoding pass and requires neither sampling nor a learned router. We study
\tailleak{} in \method{}, a dual-path open-domain question-answering framework
that applies the score to sub-question retrieval and preserves both direct and
decomposition candidates. Across four development-set traces with
Llama-3-8B-Instruct, \tailleak{} reaches \(0.749\) macro AUROC for detecting
wrong direct answers, compared with \(0.726\) for maximum token entropy. In the
tested end-to-end configuration, replacing entropy with \tailleak{} raises
macro EM/F1 from \(40.52/47.61\) to \(41.88/48.90\). Decomposition-path queries
also increase answer-containing retrieval recall on three multi-hop datasets,
while leaving single-hop PopQA nearly unchanged. These results support
\tailleak{} as a useful ranking and gating signal, but confidently wrong
answers remain a fundamental limitation.
\end{abstract}
